%%%%%%%%%%%%%%%%%%%%%%%%%%%%%%%%%%%%%%%%%%%%%%%%%%%%%%%%%%%%
%% AUFGABE 1
%%%%%%%%%%%%%%%%%%%%%%%%%%%%%%%%%%%%%%%%%%%%%%%%%%%%%%%%%%%%

\newpage
\section{Ampelsteuerung}

TODO
\begin{itemize}
    \item Aufgabenbeschreibung
\end{itemize}

\subsection{Vorbereitung}

\begin{wrapfigure}[13]{r}{6cm}
\vspace*{-11pt}
\centering
\begin{circuitikz}[scale=0.9, transform shape]
    \draw (0,0) node[ground]{}
    to[R,l^=$R_V$, i<_=\SI{15}{\milli\ampere}] ++(0,1.5)
    to[stroke led, invert, l^={LED}, v<=\SI{2}{\volt}] ++(0,1)
    |- ++(-1,0.5) node[ocirc]{} node[above right]{IO-Pin (\SI{5}{\volt})};
\end{circuitikz}
\caption{LED mit Vorwiderstand}
\end{wrapfigure}

Im Betrieb weisen die verwendeten Leuchtdioden eine Vorwärtsspannung von ca. $U_F = \SI{2}{\volt}$ auf. Um den Strom durch die Leuchtdiode zu begrenzen, wird ein Vorwiderstand so dimensioniert, dass die restliche Spannung mit dem gewählten Maximalstorm an ihm abfällt.
\[
R_V = \frac{U_\text{IO} - U_F}{I_D} = \frac{\SI{5}{\volt} - \SI{2}{\volt}}{\SI{15}{\milli\ampere}} = \SI{200}{\ohm}
\]

Zur vefügung steht die E6 Widerstandsbaureihe. Damit keinesfalls der Maximalstorm überschritten wird, wird der nächst höhere Wert aus der Reihe gewählt.
\[
R_V = \SI{220}{\ohm}
\]